\documentclass{../../../Styles/MS-Style}
\addbibresource{../../../Assets/Citations/references.bib}

\usepackage[
    pdftitle={André Mossi CV (EN)},
    pdfauthor={André Mossi},
    pdfcreator={LaTeX and NeoVim}
]{hyperref}

\begin{document}

    \begin{center}
        \authorTitle{Roberto (André) Mossi Milla}

        \contactBlock{
            Erie, Pennsylvania, 16506
        }{
            814-790-1591
        }{
            mossiroberto0392@gmail.com
        }{
            andremossi-linktree.vercel.app
        }{
            github.com/AndreM222
        }
    \end{center}

    \begin{Form}

        \sectionTitle{SUMMARY}

        \normalsize{
            Software Engineer experienced in artificial intelligence, full-stack development,
            computer vision, distributed systems, and real-time simulation environments.
            Published researcher with experience spanning academic research, industrial
            computer vision systems, AI-powered healthcare platforms, multimodal AI
            workflows, entrepreneurship, and cross-platform software architecture.
        }

        \sectionTitle{EDUCATION}

        \subSectionTitle{Gannon University}{Pennsylvania, United States}
        \educationInfo{Bachelor of Science in Computer Systems Engineering}{May 2025\cite{GannonBSDiploma}}

        \subSectionTitle{ISI Japanese Language School}{Tokyo, Japan}
        \educationInfo{Intensive Japanese Language Program}{October 2025 - March 2026\cite{ISIJapaneseSchool}}

        \subSectionTitle{Harvard University}{Online}
        \educationInfo{CS50’s Introduction to Computer Science}{December 2022\cite{HarvardCS50Certification}}

        \subSectionTitle{Extreme Networks}{Online}
        \educationInfo{Introduction to Future Networks}{May 2022\cite{ExtremeCertification}}

        \sectionTitle{RESEARCH EXPERIENCE}

        \subSectionTitle{Gannon University}{Pennsylvania, United States}

        \subSubSectionTitle{Research Assistant}{August 2024 - May 2025}

        \begin{listBlock}

            \generalItem
            Conducted research under faculty supervision on multi-agent artificial intelligence
            simulating natural selection and survival behaviors in dynamic environments.\cite{TagResearchRecommendation}

            \generalItem
            Implemented NeuroEvolution of Augmenting Topologies (NEAT) to evolve neural network architectures
            and behavioral strategies across generations.

            \generalItem
            Developed a 3D multi-agent simulation environment in Unreal Engine featuring
            procedural world generation, physics-based interactions, and reinforcement learning pipelines
            using C++, CUDA, and PyTorch.

            \generalItem
            Research paper accepted and published by the American Society for Engineering Education
            (ASEE) and presented at a Sigma Xi conference at Penn State University.\cite{AITagResearchPublication}
            \cite{AITagSigmaXIPresentation}

        \end{listBlock}

        \sectionTitle{INDUSTRY EXPERIENCE}

        \subSectionTitle{Ingeniería Digital}{Tegucigalpa, Honduras}

        \subSubSectionTitle{Full-Stack Developer / AI Systems Engineer}{May 2026 - Present}

        \textbf{AI-Assisted Medical Document Generation Platform}

        \begin{listBlock}

            \generalItem
            Developed an AI-assisted medical document generation platform for post-surgery reports (IDR),
            combining OCR extraction pipelines, LLM-powered grammar correction, and automated professional
            document structuring using Next.js, React, TypeScript, PostgreSQL, AWS S3,
            Amazon Textract, and OpenAI APIs.

            \generalItem
            Architected authenticated document management workflows using Auth0, enabling users
            to generate, analyze, edit, version, and manage AI-generated reports through
            a full-stack web platform.

            \generalItem
            Designed scalable file-processing pipelines integrating OCR extraction,
            cloud storage, AWS S3 buckets, and PostgreSQL persistence for intelligent
            document retrieval and processing.

            \generalItem
            Implemented internationalization workflows and automated translation asset generation
            using CMake to support multilingual deployments across US and Honduran clients.

        \end{listBlock}

        \textbf{AI Vehicle Damage Assessment Platform}

        \begin{listBlock}

            \generalItem
            Developed an AI-assisted vehicle damage assessment platform utilizing multimodal AI workflows,
            VIN history analysis, image-based damage detection, and insurance estimation pipelines.

            \generalItem
            Integrated computer vision and LLM APIs to analyze uploaded vehicle imagery,
            validate VIN-related information, compare against historical vehicle records,
            and support insurance decision workflows.

            \generalItem
            Developed data acquisition services using Playwright to collect and process
            vehicle information required for valuation and insurance analysis workflows.

            \generalItem
            Implemented interactive damage annotation and review tools using Konva,
            allowing users and insurance personnel to identify, visualize, and validate
            detected vehicle damage regions.

            \generalItem
            Developed responsive client-facing and administrative dashboards using React,
            Next.js, TypeScript, PostgreSQL, AWS S3, and Shadcn/UI.

            \generalItem
            Maintained engineering standards and CI/CD workflows using GitHub Actions,
            Husky, ESLint, Prettier, lint-staged, and Commitizen.

        \end{listBlock}

        \subSectionTitle{Gannon University}{Pennsylvania, United States}

        \subSubSectionTitle{Computer Vision Engineer (Project-Based) Industry Partner (NDA)}{August 2024 - May 2025}

        \begin{listBlock}

            \generalItem
            Architected and deployed a real-time industrial traceability system using Ultralytics YOLO
            to detect and track pallets, automatically generating structured events with images
            and detection metadata.\cite{ComputerVisionVisual}

            \generalItem
            Collaborated in a cross-functional team contributing across hardware, backend, and software components;
            took primary ownership of the monitoring dashboard and visualization interface.

            \generalItem
            Developed a full-stack web dashboard integrated with backend services and TimescaleDB
            to visualize detection events, images, and operational analytics in real time.

            \generalItem
            Implemented administrative tools used to curate datasets, correct detections,
            and improve model retraining workflows.

            \generalItem
            Containerized distributed services using Docker and integrated multi-language components
            (C++, Python, Rust) to coordinate inference, data processing, and communication.

        \end{listBlock}

        \subSectionTitle{Dinant}{Tegucigalpa, Honduras}

        \subSubSectionTitle{Full-Stack Developer}{June - July 2023}

        \begin{listBlock}

            \generalItem
            Led the end-to-end design and development of two production web platforms,
            conducting requirements gathering, architecture design, client presentations,
            and deployment.\cite{DinantRecommendation}

            \generalItem
            Architected a contract lifecycle management platform supporting device sales,
            rentals, repairs, and multi-country operations.

            \generalItem
            Developed a multi-role web application using Oracle APEX, JavaScript, HTML,
            and CSS supporting digital signatures, PDF generation, and administrative reporting.

            \generalItem
            Built a factory maintenance management platform enabling machine monitoring,
            maintenance workflows, repair tracking, inventory reporting, and procurement insights.

        \end{listBlock}

        \sectionTitle{PROJECTS \& ENTREPRENEURSHIP}

        \subSectionTitle{NOIR}{Founder, Product Designer \& Software Engineer — 2025 - Present}

        \begin{listBlock}

            \generalItem
            Founded and lead development of NOIR, a cinema-focused software platform
            designed to unify movie discovery, streaming availability, playback continuity,
            and cross-device viewing experiences.

            \generalItem
            Authored a comprehensive technical strategy and product design document
            defining system architecture, monetization strategy, platform integrations,
            compliance requirements, and phased market expansion.

            \generalItem
            Architected a distributed backend platform using Rust, Axum, PostgreSQL,
            Redis, Supabase, Docker, and Railway to support real-time synchronization,
            authentication, movie metadata aggregation, and availability services.

            \generalItem
            Developed native SwiftUI applications targeting macOS, iOS, iPadOS,
            and watchOS with realtime synchronization, push notifications,
            subscription management, and AI-assisted discovery workflows.

            \generalItem
            Designed engineering workflows utilizing GitHub Actions, Linear,
            Commitizen, SwiftLint, Clippy, Docker, and automated CI/CD pipelines.

        \end{listBlock}

        \subSectionTitle{Copilot-Lualine}{Open Source Project}

        \begin{listBlock}

            \generalItem
            Designed and developed Copilot-Lualine, an open-source Neovim plugin
            integrating GitHub Copilot status information into the Lualine statusline ecosystem. It holds
            as of now 116 stars and going up.

            \generalItem
            Implemented Lua-based editor tooling focused on workflow visibility,
            developer productivity, and Neovim customization.

        \end{listBlock}

        \sectionTitle{LEADERSHIP \& ORGANIZATIONS}

        \begin{listBlock}

            \generalItem
            Founder, Product Designer \& Software Engineer,
            NOIR (2025 -- Present).

            \generalItem
            President, Computer Club,
            DelCampo International School
            (2019 -- 2020).

            \generalItem
            Member, Mathematics Club, Gannon University (2023 - 2025).

        \end{listBlock}

        \sectionTitle{PUBLICATIONS}

        \begin{listBlock}

            \generalItem
            Mossi, R. A. — “Tag AI-Sandbox,” Proceedings of the ASEE Annual Conference,
            2025 (peer-reviewed).\cite{AITagResearchPublication}

        \end{listBlock}

        \sectionTitle{TECHNICAL SKILLS}

        \begin{listBlock}

            \generalItem
            Languages: C++, Rust, Swift, TypeScript, JavaScript, Python, Java,
            SQL, Lua, Bash, PowerShell.

            \generalItem
            Frontend: SwiftUI, React, Next.js, Vite, Tailwind CSS, Shadcn/UI.

            \generalItem
            Backend \& Cloud: Axum, Tokio, PostgreSQL, Redis, Supabase,
            AWS S3, Amazon Textract, Docker, Auth0, Railway, REST APIs.

            \generalItem
            AI \& Computer Vision: OpenAI APIs, Claude APIs, YOLO, CUDA,
            PyTorch, NEAT, Reinforcement Learning, OCR Pipelines,
            Multimodal AI Workflows.

            \generalItem
            Developer Tooling: GitHub Actions, Husky, ESLint, Prettier,
            Commitizen, CMake, SwiftLint, Clippy, Git.

            \generalItem
            Operating Systems: Linux, macOS, Windows.

            \generalItem
            Languages: Spanish, English,
            Japanese.

        \end{listBlock}

        \sectionTitle{AWARDS AND HONORS}

        \begin{listBlock}

            \awardItem
            {School Development Award}
            {Award obtained for outstanding development of a videogame using Blueprints in Unreal Engine 4}
            {2020\cite{SoftwareEngineer2020Award}}

            \awardItem
            {School Development Award}
            {Award obtained for outstanding development of a videogame using C++ and Unreal Engine 4}
            {2016\cite{SoftwareEngineer2016Award}}

        \end{listBlock}

        \bibliographyBlock

    \end{Form}
\end{document}
